\documentclass[12pt,a4paper,psamsfonts]{amsart}
\usepackage{amssymb,mathtools}
\usepackage{calc}
\setlength{\topmargin}{0cm}
\setlength{\oddsidemargin}{0cm}
\setlength{\evensidemargin}{0cm}
\setlength{\marginparwidth}{0cm}
\setlength{\marginparsep}{0cm}
\setlength{\textheight}{\paperheight - 2in - 35pt}
\setlength{\textwidth}{\paperwidth - 2in}
\setlength{\headheight}{12.5pt}
\setlength{\headsep}{25pt}
\setlength{\footskip}{30pt}
\renewcommand{\baselinestretch}{1.3}
\pagestyle{headings}
\raggedbottom
\setlength{\emergencystretch}{2em}
\usepackage[colorlinks=true,urlcolor=blue,linkcolor=blue,citecolor=blue]{hyperref}
\newtheorem*{theorem}{Theorem}
\theoremstyle{definition}
\newtheorem*{setup}{Setup}
\theoremstyle{plain}
\newtheorem{lemma}{Lemma}[section]
\newtheorem{proposition}[lemma]{Proposition}
\newtheorem{corollary}[lemma]{Corollary}
\newcommand{\Q}{\mathbb Q}
\newcommand{\C}{\mathbb C}
\newcommand{\OO}{\mathcal O}
\newcommand{\mult}{\operatorname{mult}}
\newcommand{\nef}{\operatorname{n}}
\title{Nef Abundance: Framework}
\date{Temporary draft -- October 11, 2026}
\begin{document}
\maketitle

\section{Theorem and proof}

\begin{theorem}
Let $X$ be a smooth projective threefold over $\C$ such that
$K_X\sim0$ and $q(X)=0$. Let $D$ be a nef Cartier divisor on $X$
such that $\nef(X,D)=3$ and $D^3=0$. Then
\[
 H^0(X,\OO_X(mD))\ne0
 \qquad\text{for some integer }m>0.
\]
\end{theorem}

\begin{proof}
Suppose for contradiction that $\kappa(D)=-\infty$, equivalently
$H^0(X,\OO_X(mD))=0$ for every integer $m>0$.

\begin{setup}
\leavevmode
\begin{itemize}
\setlength{\itemsep}{2pt}
\setlength{\parsep}{0pt}
\item Fix an ample Cartier divisor $A$ and $\epsilon\in\Q_{>0}$.
For integers $n\ge1$, put
\[
 a_n=((D+A/n)^3)^{1/3},\qquad
 r_n=\max\{1,\lfloor\epsilon/a_n+1/2\rfloor\},\qquad
 P_n=r_n(D+A/n).
\]
Then $P_n$ is ample and rational Cartier, $r_n\to\infty$, and
$P_n^3\to\epsilon^3$. Indeed, $a_n\to0$ and
$|r_na_n-\epsilon|\le a_n/2$ for all sufficiently large $n$.
\item Let $p_i:X\times X\to X$ be the projections, for $i=1,2$.
Put
\[
 Z=\mathbb P_{X\times X}
       \bigl(p_1^*\OO_X(D)\oplus p_2^*\OO_X(D)\bigr).
\]
Let $S_1,S_2\subset Z$ be the two summand sections and set
\[
 \mathbb T=Z\setminus(S_1\cup S_2).
\]
\item Let $D_i,P_{n,i}$ denote the pullbacks of $D,P_n$ by $p_i$,
and also their pullbacks to $Z$.
Set $\xi=c_1(\OO_Z(1))$. For an integer $N>0$, put
\[
 R_n=P_{n,1}+P_{n,2}+N\xi.
\]
Whenever $|kR_n|$ is used, $k>0$ is an integer with $kP_n$ Cartier.
\end{itemize}

For $x\in X$, let $\mathcal I_x$ be the ideal sheaf of the point $x$.
An \emph{exceptional functional} is a linear functional
$\Phi:N^1(X)_{\mathbb R}\to\mathbb R$ supplied with data
$(x,\pi:V\to X,H_1,H_2)$, where $V$ is a smooth projective threefold,
$\pi$ is birational, the extended ideal
\[
 \mathcal I_x\OO_V:=\operatorname{Im}(\pi^*\mathcal I_x\to\OO_V)
\]
is invertible and equals $\OO_V(-E_V)$ for an effective
$\pi$-exceptional Cartier divisor $E_V$, and $H_1,H_2$ are ample
Cartier divisors on $V$. Define
\[
 \Phi(L)=\frac{\pi^*L\cdot H_1\cdot H_2}{E_V\cdot H_1\cdot H_2}.
\]
All intersections are on $V$, and the denominator is positive since
$E_V\ne0$. 
For every effective Cartier divisor $F$ on $X$, $\mult_x(F)\le\Phi(F)$:
a local equation of $F$ is a section of $\mathcal I_x^{\mult_x(F)}$
near $x$, so $\pi^*F\ge\mult_x(F)E_V$; intersect with $H_1,H_2$
and divide by the positive denominator.

For a Cartier divisor $L$ on a smooth projective variety $W$,
$|L|$ \emph{separates $s$-jets at $z$} if
\[
 \begin{gathered}
 \forall u\in\OO_{W,z}(L),\quad
 \exists\sigma\in H^0(W,\OO_W(L)):\\
 \sigma_z-u\in\mathfrak m_z^{s+1}\OO_{W,z}(L),
 \end{gathered}
\]
where $s\ge0$ is an integer and $\sigma_z$ denotes the germ at $z$.
We take $\epsilon$ sufficiently small that
\[
 343\epsilon^3<\tfrac14,\qquad40\epsilon<\tfrac38.
\]
\end{setup}

We use the following three lemmas under the above setting.
Their proofs are not included in this temporary framework.

\begin{lemma}\label{lem:framework-witness}
Let $X$ be a smooth projective threefold over $\C$ with $K_X\sim0$
and $q(X)=0$. Let $D$ be a nef Cartier divisor with $\nef(X,D)=3$,
$D^3=0$ and $\kappa(D)=-\infty$, and let $A$ be an ample Cartier
 divisor. Fix $\epsilon\in\Q_{>0}$ with $343\epsilon^3<1/4$ and
$40\epsilon<3/8$. Let $P_n=r_n(D+A/n)$ be the sequence defined in Setup.
For every integer $N>0$, there exists an integer
$n_0=n_0(\epsilon,N)\ge1$ such that, for every integer $n\ge n_0$,
there exists an exceptional functional $\Phi$ with
$\Phi(P_n+ND)<40\epsilon$.
The functional and its presentation may depend on $n,N$;
$n_0$ is independent of these subsequent choices.
\end{lemma}

\begin{lemma}\label{lem:framework-obstruction}
Let $X$ be a smooth projective threefold over $\C$ with $K_X\sim0$
and $q(X)=0$. Let $D$ be a nef Cartier divisor with $\nef(X,D)=3$
and $D^3=0$, and let $A$ be an ample Cartier divisor.
Fix $\epsilon\in\Q_{>0}$ with $343\epsilon^3<1/4$, and let $P_n,R_n$
be defined as in Setup.
For every integer $N>0$, there exists an integer
$n_1=n_1(\epsilon,N)\ge1$ such that, for every integer $n\ge n_1$,
if there exists an exceptional functional $\Phi$ with
$\Phi(P_n+ND)<\tfrac38$, then $|kR_n|$ does not
separate $8k$-jets at any $z\in\mathbb T$, for any integer $k>0$
with $kP_n$ Cartier.
\end{lemma}
For fixed $N$, take $n_1$ so that $343P_n^3<1/4$ and $r_n>N$ for
all $n\ge n_1$, using $P_n^3\to\epsilon^3$, $343\epsilon^3<1/4$
and $r_n\to\infty$. Since $P_n-r_nD$ is nef,
\[
 (2P_n+ND)\cdot P_n^2\le(2+N/r_n)P_n^3<4P_n^3.
\]
These are the numerical conditions for the geometric input in
Lemma~\ref{lem:framework-obstruction}; its proof is still not included.

\begin{lemma}\label{lem:framework-jets}
Let $X$ be a smooth projective threefold over $\C$ with $K_X\sim0$
and $q(X)=0$. Let $D$ be a nef Cartier divisor with $\nef(X,D)=3$
and $D^3=0$, and let $A$ be an ample Cartier divisor.
Assume $\kappa(D)=-\infty$. Fix $\epsilon\in\Q_{>0}$ with
$343\epsilon^3<1/4$ and $40\epsilon<3/8$, and let $P_n,Z,R_n$
be defined as in Setup.
There exists an integer $N>0$ such that, for every $n\gg0$,
$|kR_n|$ separates $8k$-jets at some $z\in\mathbb T$ for some
integer $k>0$ with $kP_n$ Cartier.
\end{lemma}

\newpage
We now complete the proof under the standing assumption
$\kappa(D)=-\infty$. Lemma~\ref{lem:framework-witness} uses this
assumption to produce a small exceptional functional;
Lemma~\ref{lem:framework-obstruction} uses that functional to prohibit
jet separation, without requiring $\kappa(D)=-\infty$.
Lemma~\ref{lem:framework-jets} uses $\kappa(D)=-\infty$ to provide
jet generation. Thus these geometric inputs rule out
$\kappa(D)=-\infty$ under the common geometric hypotheses.
Choose $N$ by Lemma~\ref{lem:framework-jets}. For this fixed $N$, take
$n\ge\max\{n_0(\epsilon,N),n_1(\epsilon,N)\}$ also large enough for
Lemma~\ref{lem:framework-jets}. By Lemma~\ref{lem:framework-witness},
there exists $\Phi$ such that
\[
 \Phi(P_n+ND)<40\epsilon<\tfrac38.
\]
Lemma~\ref{lem:framework-obstruction} then says that $|kR_n|$ does not
separate $8k$-jets at any $z\in\mathbb T$ for any integer $k>0$ with
$kP_n$ Cartier. But Lemma~\ref{lem:framework-jets}, for the same $N,n$,
provides such an integer $k$ and a point $z\in\mathbb T$ at which
$|kR_n|$ separates $8k$-jets. This is a contradiction. Hence
$H^0(X,\OO_X(mD))\ne0$ for some integer $m>0$.\qedhere
\end{proof}

\medskip
\noindent\textbf{Draft status.}
This is a reformulation of the first-section argument with its three
geometric inputs retained, rather than an independent verification
of those inputs. The chosen presentation of $\Phi$, including the
model $\pi:V\to X$, the divisors $E_V,H_1,H_2$, and the extended-ideal identity
$\mathcal I_x\OO_V=\OO_V(-E_V)$ must be fixed before reduction of
characteristic; the bare linear functional does not supply these data. This AI-assisted
draft is for discussion and checking. No claim of novelty,
priority, or authorship credit is made.

\end{document}
